% Source: https://artofproblemsolving.com/community/c3h1770260_geometry
% Stable identity: aops-topic-1770260-q01
Rotate triangle $BPC$ about $B$ through $60^\circ$, in the direction that takes $C$ to $A$ (possible because $BC=BA$ and $\angle CBA=60^\circ$). Let $P$ go to $P'$.

\begin{asy}
usepackage("amsmath");
size(5.5cm);
defaultpen(fontsize(10pt));
real s3=sqrt(3);
pair B=(0,0),C=(4,0),A=(2,2*s3);
pair P=(2,-2*s3)+4*dir(80);
pair Pp=rotate(60,B)*P;
draw(A--B--C--cycle,linewidth(0.9));
draw(P--A);draw(P--B);draw(P--C);
draw(B--Pp,blue+dashed);draw(P--Pp,blue+dashed);draw(A--Pp,blue+dashed);
pair u=unit(A-Pp),v=unit(P-Pp);
draw(Pp+0.25*u--Pp+0.25*u+0.25*v--Pp+0.25*v,blue);
dot(P);dot(Pp,blue);
label("$A$",A,N);label("$B$",B,SW);label("$C$",C,SE);
label("$P$",P,dir(25));label("$P'$",Pp,W,blue);
label("\mbox{}",(min(currentpicture,user=true)+max(currentpicture,user=true))/2,black); // keeps the text after the figure black when LaTeX typesets the labels inline
\end{asy}

Then $BP'=BP$ and $\angle PBP'=60^\circ$, so triangle $BPP'$ is equilateral: $PP'=BP$ and $\angle BP'P=60^\circ$. Also $AP'=CP$ and $\angle BP'A=\angle BPC$, since the rotation carries triangle $BPC$ onto triangle $BP'A$.

In triangle $AP'P$,
\[
AP'^2+PP'^2=CP^2+BP^2=AP^2 ,
\]
so by the converse of the Pythagorean theorem $\angle AP'P=90^\circ$. The point $P$ lies inside the triangle and $P'$ outside it beyond side $AB$, so the ray $P'B$ lies on the far side of $P'P$ from $A$ and
\[
\angle BPC=\angle BP'A=\angle BP'P+\angle PP'A=60^\circ+90^\circ=150^\circ .\qquad\blacksquare
\]
